% ENHANCEMENT.tex, the manuscript skeleton of the enhancement paper of
% papers/wealth_garch. Created 25 October 2026 as pass 1 of the plan in
% enhancement/TODO.md, following papers/MANUSCRIPT_SKELETON.md, with the parent
% MANUSCRIPT.tex as the reference for the empirical rows. The author asked
% for the write-up on 25 October 2026, after the 2025 hold-out.
%
% STRUCTURE THE CHECKER RELIES ON (enhancement/checks/verify_enhancement.py)
%   Rows are written only through the seven row macros below, and proofs only
%   through the mproof environment. IDs are K1, ... (introduction), P1, ...
%   (players), A1, ... (assumptions), M1, ... (model), E1, ... (empirical),
%   L1, ... (literature) and C1, ... (conclusions).
%     \krow{ID}{headline}{model or empirical rows}{scope note}
%     \prow{ID}{player}{choice faced}{how the model settles it}{in the literature}{rows}
%     \arow{ID}{assumption}{in plain English}{symbols}{model rows}{reason}
%     \mrow{ID}{claim in words}{statement in symbols}{Lean theorems}
%     \erow{ID}{claim}{verdict}{model rows}{map rows}{evidence files}
%     \lrow{ID}{bib key}{what we rely on}{verbatim quote}{page}{rows}
%     \crow{ID}{statement}{model, literature or empirical rows}
%     \begin{mproof}{model row}{Lean theorems} ... \end{mproof}
%   Lean names are fully qualified names declared under lean/mathlib/. Map rows
%   are the G rows of enhancement/MEASUREMENT_MAP.tex. Evidence paths are
%   relative to papers/wealth_garch.
%
% RULES THE CHECKER ENFORCES
%   The parent rules MS-1 to MS-7 (checks/verify_manuscript.py) with this
%   file's map, and MS-8 (assumption rows point to model rows, and every model
%   row rests on an assumption) and MS-9 (player rows point to model,
%   empirical or literature rows). Every empirical row's results files are
%   re-run on their pinned data and must reproduce exactly, each committed
%   after the pre-registrations it names.
\documentclass[11pt]{article}
\usepackage[T1]{fontenc}
\usepackage[utf8]{inputenc}
\usepackage{mathptmx}
\usepackage{microtype}
\usepackage[a4paper,margin=1.8cm]{geometry}
\usepackage{amsmath,amssymb,amsthm}
\usepackage{booktabs}
\usepackage{array}
\usepackage{longtable}
\usepackage{url}
\usepackage[hidelinks]{hyperref}
\usepackage{natbib}

\emergencystretch=1.5em
\DeclareUnicodeCharacter{2019}{'}

\newcolumntype{P}[1]{>{\raggedright\arraybackslash}p{#1}}

\urlstyle{tt}
\makeatletter
\g@addto@macro\UrlBreaks{\do\.\do\_\do\,\do\/}
\makeatother
\newcommand{\lean}[1]{\path{#1}}
\newcommand{\rowref}[1]{#1}
\newcommand{\unv}[1]{\textit{Unverified:} #1}
\newcommand{\pend}[1]{\textit{Pending:} #1}

\newcommand{\krow}[4]{#1 & #2 & \rowref{#3} & #4 \\ \addlinespace}
\newcommand{\prow}[6]{#1 & #2 & #3 & #4 & #5 & \rowref{#6} \\ \addlinespace}
\newcommand{\arow}[6]{#1 & #2 & #3 & #4 & \rowref{#5} & #6 \\ \addlinespace}
\newcommand{\mrow}[4]{#1 & #2 & #3 & \lean{#4} & p.~\pageref{proof:#1} \\ \addlinespace}
\newcommand{\erow}[6]{#1 & #2 & #3 & \rowref{#4} & \rowref{#5} & \path{#6} \\ \addlinespace}
\newcommand{\lrow}[6]{#1 & \citet{#2} & #3 & ``#4'' & #5 & \rowref{#6} \\ \addlinespace}
\newcommand{\crow}[3]{#1 & #2 & \rowref{#3} \\ \addlinespace}
\newcommand{\mm}[6]{#1 & \textbf{#2} & #3 & #4 & #5 & #6 \\ \addlinespace}

\newenvironment{mproof}[2]
  {\subsection*{Proof of #1}\label{proof:#1}%
   {\raggedright\noindent\textit{Lean:} \lean{#2}.\par}\smallskip}
  {\par\hfill$\square$\par\medskip}

\newcommand{\fl}{\underline{k}}

\title{Bank Capital Buffers and the Floor They Guard}
\date{Skeleton of 25 October 2026}

\begin{document}
\maketitle

\section{Introduction}\label{sec:intro}

\subsection{Question}

Why does the standard dynamic model of bank capital put the buffers that US
banks hold far below what they hold, and does a model that guards the floor
regulation makes binding, with the risk that book capital carries, place them
better?

\subsection{Answer}

The standard model, the dividend barrier of a bank whose capital moves with
its own earnings and which fails at the legal minimum, predicts that banks pay
dividends from thin buffers and almost never hold capital far above the
reference, and observed banks do the opposite in every window scored
(E1). The buffers banks hold match that model only when the risk they guard
against is about three times their earnings volatility (E2). Measures built
from losses, the supervisor's stress-test decline and the bank's own worst
decline, put that risk at about one times earnings volatility, and the
market's equity volatility at about eight times (E3).

A barrier model whose floor is the line the supervisor's restrictions start
at, the well-capitalised ratio before 2010 and the minimum plus the
conservation buffer after, and whose risk is the observed movement of book
capital, beats the standard model on banks it had not been scored on (E4,
E5), and on 2025, a year kept back from every choice (E6). It does not place
the level of buffers on 2025, where it puts too few bank-quarters far above
the reference (E7). The floor enters the barrier model only as a shift
of the buffer (M3), so the floor changes where the barrier sits in the ratio
and not the barrier itself (M2).

Two facts about buffers stay outside both models. Banks damp the movement of
their own capital ratios by about a sixth through their capital actions (E8).
Smaller banks hold larger buffers at the same risk, and neither access to a
parent nor a listing accounts for it (E9).

\begin{center}\footnotesize
\begin{longtable}{P{0.5cm}P{6.0cm}P{1.6cm}P{7.6cm}}
\caption{Headlines.}\label{tab:headlines}\\
\toprule ID & Headline & Rows & Scope note \\ \midrule \endfirsthead
\toprule ID & Headline & Rows & Scope note \\ \midrule \endhead
\krow{K1}{The standard dividend-barrier model of bank capital fails both facts in every window, on US commercial banks 2005 to 2025.}{E1, E6, M2}{Scored on the two pre-registered facts only (the share of dividend quarters paid from a buffer below one earnings volatility over the reference, and the share of bank-quarters more than four above it). Peura and Keppo report a shortfall in level themselves (L1).}
\krow{K2}{Buffers imply a risk about three times earnings volatility, between the one times of loss-based measures and the eight times of the market.}{E2, E3}{Training window only, in the non-issuer barrier model. Risk scales, not mean capital shortfalls under a fixed scenario (L13). The implied volatility is Peura and Keppo's idea, at about 1.1 to 1.3 times in their sample (L15).}
\krow{K3}{A barrier model guarding the regulatory floor with book-capital risk beats the standard model out of sample and on a hold-out year.}{E4, E5, E6, M3}{Beats means a smaller sum of the two misses. The floor was set from sources after a floor near the reference had been seen to work (E5, E6). It does not place the level of buffers on the hold-out (E7).}
\krow{K4}{Banks' own capital actions damp the movement of their capital ratios, and smaller banks hold more capital at the same risk.}{E8, E9}{Description of observed data. The book-capital risk of K3 contains the damping, and removing it makes the model overshoot (E8).}
\bottomrule
\end{longtable}
\end{center}

\section{Model}\label{sec:model}

\subsection{Players}

\begin{center}\footnotesize
\begin{longtable}{P{0.5cm}P{1.8cm}P{2.6cm}P{3.6cm}P{4.7cm}P{1.4cm}}
\caption{Players.}\label{tab:players}\\
\toprule ID & Player & Choice faced & How the model settles it & In the literature & Rows \\ \midrule \endfirsthead
\toprule ID & Player & Choice faced & How the model settles it & In the literature & Rows \\ \midrule \endhead
\prow{P1}{Bank}{When to pay dividends.}{Derived. The bank pays at the dividend barrier, the point of smooth fit (M2).}{Peura and Keppo's non-issuer bank, the case Milne and Whalley's buffer reduces to (L1, L3).}{M2, L1, L3}
\prow{P2}{Supervisor}{Where to set the floor and when to restrict a bank.}{Taken as given. The floor is a rule (A5), with its post-2009 value from the conservation buffer (M5).}{Orozco and Rubio's well-capitalised threshold and its replacement after 2009 (L7, L8).}{M5, L7, L8}
\prow{P3}{Insurer}{How to price deposit insurance.}{Not modelled. The insurer is state, entering only through the floor banks guard.}{The lower assessment above the well-capitalised line is one reason banks guard it (L7).}{L7}
\prow{P4}{Investors}{How to value bank equity.}{Not modelled. Investors enter only as one measure of the risk the bank guards against (E3).}{Craine and Martin find market capital moves far more than book capital (L12).}{E3, L12}
\bottomrule
\end{longtable}
\end{center}

\subsection{Assumptions}

\begin{center}\footnotesize
\begin{longtable}{P{0.5cm}P{3.6cm}P{4.2cm}P{1.8cm}P{1.0cm}P{3.6cm}}
\caption{Assumptions.}\label{tab:assumptions}\\
\toprule ID & Assumption & In plain English & Symbols & Model rows & Reason \\ \midrule \endfirsthead
\toprule ID & Assumption & In plain English & Symbols & Model rows & Reason \\ \midrule \endhead
\arow{A1}{Between dividend payments $dx_t=\mu_i\,dt+s_i\,dW_t$.}{The buffer over the floor drifts up with the bank's mean earnings and moves with the judged risk through a Brownian motion.}{$x$, $\mu_i$, $s_i$, $W_t$}{M1, M2}{Peura and Keppo's capital process (Appendix D, process).}
\arow{A2}{Dividends are paid only at the barrier $u_0$, and no equity is issued.}{The bank pays out at the dividend barrier and never raises capital.}{$u_0$}{M2}{The non-issuer case, an upper bound on the buffer when issuance is costly (L3; Appendix D, issuance).}
\arow{A3}{The bank is closed when $x_t=0$.}{The bank fails when the buffer over the floor is used up.}{$x$}{M1, M2, M3}{Closure at the floor (Appendix D, closure).}
\arow{A4}{Dividends are discounted at $\rho=0.04$.}{The bank values dividends at a discount rate of four percent a year, its cost of equity.}{$\rho$}{M1}{Peura and Keppo's value (Appendix D, discounting).}
\arow{A5}{$\fl=0.08$ in the baseline, and $\fl=0.10$ to 2009Q4 and $\fl=0.105$ from 2010Q1 in the enhanced model.}{The regulatory floor is the legal minimum in the baseline, and the well-capitalised line and then the minimum plus the conservation buffer in the enhanced model.}{$\fl$}{M3, M5}{The lines the sources find binding (L7, L8, L9; Appendix D, floor).}
\arow{A6}{$s_i=\sigma_i$ in the baseline and $s_i=\sigma_{D,i}$ in the enhanced model.}{The judged risk is the earnings volatility in the baseline and the book-capital volatility in the enhanced model.}{$s_i$, $\sigma_i$, $\sigma_{D,i}$}{M1, M2, M4}{Appendix D, judged risk.}
\bottomrule
\end{longtable}
\end{center}

\subsection{Results}

\begin{center}\footnotesize
\begin{longtable}{P{0.5cm}P{4.2cm}P{5.0cm}P{4.2cm}P{0.9cm}}
\caption{Model results.}\label{tab:model}\\
\toprule ID & Claim & Statement & Lean & Proof \\ \midrule \endfirsthead
\toprule ID & Claim & Statement & Lean & Proof \\ \midrule \endhead
\mrow{M1}{The value equation has one positive and one negative root.}{$s\neq0$, $\rho>0$ $\Rightarrow$ $p=d_+(\mu,s,\rho)>0>m=d_-(\mu,s,\rho)$, both roots of $\tfrac12 s^2d^2+\mu d-\rho=0$.}{PeuraKeppo.char_eq, PeuraKeppo.dplus_pos, PeuraKeppo.dminus_neg}
\mrow{M2}{The dividend barrier satisfies smooth fit and equals Milne and Whalley's buffer.}{$p>0>m$ $\Rightarrow$ $a_1e^{-pu_0}+a_2e^{-mu_0}=0$ at $u_0=\frac{2}{p-m}\log\frac{-m}{p}$, with $a_1>0>a_2$, and $\frac{1}{p-m}\log\frac{m^2}{p^2}=u_0$.}{PeuraKeppo.no_issue_barrier, PeuraKeppo.coeff_signs, MilneWhalley.buffer_eq_u0}
\mrow{M3}{Raising the floor shifts the buffer and leaves its movements unchanged.}{$(C_t+d)-(C_0+d)=C_t-C_0$ for every $d$.}{MilneWhalley.buffer_shift_invariant}
\mrow{M4}{A stress-test decline over nine quarters converts to a judged risk.}{$D/(z\sqrt{9/4})=D/(\tfrac32 z)$, and $m\le a\Rightarrow a-m\ge0$.}{StressTests.judged_sigma, StressTests.sqrt_horizon, StressTests.decline_of_rows}
\mrow{M5}{After 2009 the floor is the reference and lies above the well-capitalised line.}{$0.08+0.025=0.105$ and $0.10<0.105$.}{OrozcoRubio.basel3_floor_eq_R, OrozcoRubio.basel3_floor_above_pca}
\bottomrule
\end{longtable}
\end{center}

\section{Empirical analysis}\label{sec:empirical}

Every row below is a results file that the checker re-runs on its pinned
data, committed after the pre-registrations it names. F1 is the share of
dividend quarters paid from a buffer below $\sigma_i$ over $R$, and F2 the
share of bank-quarters more than $4\sigma_i$ above $R$. The
miss of a model is the sum of its two distances to the observed shares, in
percentage points.

\begin{center}\footnotesize
\begin{longtable}{P{0.5cm}P{5.2cm}P{1.0cm}P{0.9cm}P{2.2cm}P{5.0cm}}
\caption{Empirical rows.}\label{tab:empirical}\\
\toprule ID & Claim & Verdict & Model rows & Map rows & Evidence \\ \midrule \endfirsthead
\toprule ID & Claim & Verdict & Model rows & Map rows & Evidence \\ \midrule \endhead
\erow{E1}{The baseline misses both facts in both windows. On 2015 to 2024 for the 61 largest banks it predicts F1 90.0\% against 4.8\% and F2 0.0\% against 61.6\%.}{holds}{M2}{thin-buffer dividend share, far-above share}{comparison/PREREG_comparison.md, comparison/PREREG_comparison_2.md, comparison/results_pk.md}
\erow{E2}{In the non-issuer barrier model the training-window F2 of the 61 largest banks (48.0\%) is reached when the judged risk is about three times earnings volatility (42.8\% at three times, minimum 8\%).}{holds}{M2}{earnings volatility, judged risk, far-above share}{comparison/PREREG_comparison_6.md, comparison/PREREG_comparison_2.md, comparison/results_scale.md}
\erow{E3}{Stress-test declines give a judged risk with median 0.99 times earnings volatility and the bank's own worst decline 1.00 times, equity volatility scaled by the capital ratio 7.75 times.}{holds}{M4}{judged risk, market judged risk, stress-test judged risk, own-loss judged risk}{comparison/PREREG_comparison_7.md, comparison/PREREG_comparison_8.md, comparison/PREREG_comparison_9.md, comparison/PREREG_comparison_2.md, comparison/results_judged.md, comparison/PREREG_comparison_16.md, comparison/results_own_loss.md}
\erow{E4}{With equity volatility as the judged risk the barrier model beats the baseline on 2015 to 2024 on the 17 stress-tested banks (miss 30.0 against 139.8) and on 18 further banks scored blind (32.8 against 149.3).}{holds}{M2}{market judged risk, thin-buffer dividend share, far-above share}{comparison/PREREG_comparison_10.md, comparison/PREREG_comparison_7.md, comparison/PREREG_comparison_8.md, comparison/PREREG_comparison_9.md, comparison/results_test_judged.md, comparison/PREREG_comparison_12.md, comparison/PREREG_comparison_13.md, comparison/PREREG_comparison_14.md, comparison/results_extend.md}
\erow{E5}{With book-capital volatility and a 10\% floor the barrier model beats the baseline on 26 banks scored blind on 2015 to 2024 (miss 15.2 against 149.2).}{holds}{M2, M3}{regulatory floor, book-capital volatility, thin-buffer dividend share, far-above share}{comparison/PREREG_comparison_17.md, comparison/results_book_vol.md}
\erow{E6}{The frozen model, with book-capital volatility and the floor of A5, beats the baseline on the 2025 hold-out for 42 banks (miss 23.3 against 182.4) and agrees on F1 (0.2\% against 1.6\%).}{holds}{M2, M3, M5}{regulatory floor, book-capital volatility, thin-buffer dividend share, far-above share}{comparison/PREREG_comparison_20.md, comparison/results_holdout.md, comparison/results_frozen.md}
\erow{E7}{The frozen model matches the level of buffers on the 2025 hold-out. It predicts F2 67.1\% against 89.1\% observed, 22.0 points low.}{fails}{M2, M3, M5}{far-above share}{comparison/PREREG_comparison_20.md, comparison/results_holdout.md}
\erow{E8}{The capital ratio before the bank's own capital actions moves 1.19 times as much as the ratio after them (median, the 61 largest banks), and with that risk the model overshoots F2 (71.2\% against 48.0\%).}{holds}{M2}{book-capital volatility, do-nothing volatility}{comparison/PREREG_comparison_19.md, comparison/results_do_nothing.md}
\erow{E9}{Across banks of all sizes the median buffer falls with bank size (Spearman $-0.34$ and $-0.39$) and rises with risk ($+0.23$, $+0.26$). Neither model reproduces the size gradient, and stand-alone banks and banks without a listed parent hold no more than others of the same size.}{holds}{M2}{buffer over the floor, earnings volatility, bank size}{comparison/PREREG_comparison_17.md, comparison/results_book_vol.md, comparison/PREREG_comparison_18.md, comparison/results_access.md}
\bottomrule
\end{longtable}
\end{center}

\section{Literature}\label{sec:literature}

\begin{center}\footnotesize
\begin{longtable}{P{0.5cm}P{1.8cm}P{4.2cm}P{6.0cm}P{0.8cm}P{1.4cm}}
\caption{Literature read.}\label{tab:literature}\\
\toprule ID & Paper & What we rely on & Quote & Page & Rows \\ \midrule \endfirsthead
\toprule ID & Paper & What we rely on & Quote & Page & Rows \\ \midrule \endhead
\lrow{L1}{peura2006}{The baseline's own authors report it under-explains the level of capital (E1). Their fit used their own barriers, not common facts.}{The model falls somewhat short of explaining the level of capital held by banks.}{2188}{M2, C1}
\lrow{L2}{peura2006}{Their best fit, on a subsample (E1 scores all banks).}{the model explains some 40\% of the variation in capital ratios across banks and 70\% of the average level of the capital buffer.}{2191}{C1}
\lrow{L3}{milnewhalley2001}{The barrier without issuance is an upper bound on the buffer when issuance is costly (A2, M2).}{so the closed form solution represents an upper bound on the capital buffer.}{42}{M2}
\lrow{L4}{berger2008}{Targets far above the regulatory line (E1, E2).}{choose target capital levels substantially above well-capitalized regulatory minima}{0}{C1}
\lrow{L5}{berger2008}{Targets rise with risk and fall with size (E9).}{these targets increase with BHC risk but decrease with BHC size}{0}{C7}
\lrow{L6}{berger2008}{Their explanations of the size gradient, offered and not separated (E9 tests one).}{Finally, bank capital ratios tend to be inversely related to asset size. This might reflect greater diversification, scale economies in risk management, and/or lower expected cost of raising new equity on short notice.}{7}{C7}
\lrow{L7}{orozcorubio2025}{Banks guard the well-capitalised line before 2010 (A5). Their evidence, not ours.}{a significant discontinuity at the 10\% regulatory capital threshold, where banks receive benefits for exceeding it.}{1421}{M5, C3, C4}
\lrow{L8}{orozcorubio2025}{The post-2009 floor (A5, M5).}{That is, this new buffer increases the minimum reg-ulatory capital to at least 10.5\%.}{1430}{M5}
\lrow{L9}{orozcorubio2025}{The line that binds is 10\%, not 8\%, their reading of Gropp and Heider (L11).}{Our results align with the notion that capital is costly but suggest that the 10\% ratio is the relevant minimum, elucidat-ing why previous studies find weak evidence of binding capital requirements at the 8\% regulatory minimum}{1425}{M5}
\lrow{L10}{orozcorubio2025}{Their design is local to the line, so the level of buffers far above it (F2) is not theirs.}{Our main sample is restricted to banks that reported regulatory capital between 8\% and 12\%.}{1427}{C3}
\lrow{L11}{groppheider2009}{Regulation second order for large listed banks, tested at the Basel Tier 1 minimum (L9 answers).}{The paper shows that mispriced deposit insurance and capital regulation were of second order importance in determining the capital structure of large U.S. and European banks during 1991 to 2004.}{4}{C4}
\lrow{L12}{crainemartin2017}{Book capital moves little, market capital much (E3). Theirs is dispersion across banks at one date (E8 is over time).}{Book value capital ratios hardly move in response to severe economic distress.}{1}{C2}
\lrow{L13}{crainemartin2017}{The stress test and the market call for similar average changes in capital (K2 is a risk scale).}{is close to the average deviation for the market value measures (columns 5 and 6).}{14}{C2}
\lrow{L14}{langmenno2026}{A buffer requirement above the minimum binds in theory (A5). A two-period model with no buffers far above it.}{Banks do not maintain voluntary capital buffers above the CBR}{18}{M5}
\lrow{L15}{peura2006}{They introduce the implied volatility that reproduces buffers, about 1.3 times their estimate for large banks on 1983 to 2002 (C6 measures it again, larger).}{For large banks (average assets over 40 billion US\$), the average capital buffer of 3.8\% implies a volatility of 1.06\%.}{2190}{C6}
\lrow{L16}{peura2006}{The size gradient beyond volatility, which they assign to capital-market access (C7 tests that reading).}{suggesting that larger banks hold less capital than small banks, even after controlling for possible differences in their}{2190}{C7}
\lrow{L17}{daihuangkeppo2019}{Two versions of one dynamic model compared out of sample on US banks, scored on equity value, dividends and recapitalisations (C5 compares models of different authors on buffer facts).}{the partially observed model explains the out-of-sample variations substantially better than the fully observed model.}{86}{C5}
\lrow{L18}{daihuangkeppo2019}{Reported volatility understates the volatility behind it, the nearest antecedent of C6 and of the damping in E8.}{they are able to hide about one-third of the true underlying asset volatility.}{84}{C6}
\lrow{L19}{juseliusetal2025}{A buffer of about 3.8 standard deviations of buffer shocks after the crisis, for 34 large groups. The denominator is the volatility of buffer changes net of the bank's actions, not earnings volatility, and no model is asked to reproduce the buffer (C6).}{The implication is that an eight-standard-deviation shock is needed to wipe out the MB target on impact in the former regime, whereas a four-standard-deviation shock would have a similar impact in the latter regime.}{12}{C6}
\lrow{L20}{corbaederasmo2021}{A dynamic model of US banks estimated to buffer levels by size in sample, 1984 to 2007 (C5 predicts without fitting, out of sample).}{After estimating parameters to match observed levels of buffers in both the imperfect and perfect competition economies}{31}{C5}
\lrow{L21}{corbaederasmo2021}{Funding volatility as the size channel, calibrated and not tested (C7 has not tested it).}{big banks have less volatile funding inflows, which is one important factor explaining why they hold a smaller capital buffer in our model.}{24}{C7}
\lrow{L22}{senguptahogue2014}{Banks with no holding company hold more Tier 1 capital within banks, community banks 1996 to 2006. E9 pools one-bank holding companies with them and compares across banks of a size.}{Non-BHC banks are associated with a 1.06-percentage-point higher Tier 1 leverage ratio and a 2.19-percentage-point higher Tier 1 risk-based capital ratio, after controlling for other determinants of capital ratios.}{16}{C7}
\lrow{L23}{peurajokivuolle2003}{They invert a buffer model on the median US capital ratio for a confidence level, with the risk model fixed (C6 inverts for a risk scale with the decision rule fixed).}{Given a bank’s portfolio and its capital ratio, there is an implied value of a which makes the capital ratio solved from the model equal to the observed capital ratio.}{20}{C6}
\lrow{L24}{peurajokivuolle2003}{A deterministic stress scenario carries no probability, so the quantile $z$ of M4 is assumed.}{The obvious shortcoming of such an analysis is the absence of any likelihoods attached to the chosen rating transition scenario itself.}{36}{M4}
\bottomrule
\end{longtable}
\end{center}

Papers cited and not read are not used.

\section{Conclusions}\label{sec:conclusions}

Table~\ref{tab:conclusions} holds the conclusions that rest on model rows,
empirical rows and papers read in full. Table~\ref{tab:pending} holds every
claim that something is new and still waits on something, and what.

\begin{center}\footnotesize
\begin{longtable}{P{0.5cm}P{13.0cm}P{2.6cm}}
\caption{Conclusions.}\label{tab:conclusions}\\
\toprule ID & Statement & Rows \\ \midrule \endfirsthead
\toprule ID & Statement & Rows \\ \midrule \endhead
\crow{C1}{The standard dividend-barrier model of bank capital fails on US commercial banks because it guards the legal minimum with the bank's earnings volatility, and observed buffers sit where a risk three times larger would put them.}{E1, E2, L1}
\crow{C2}{Neither loss-based measure, the stress test or the bank's own worst decline, supplies that risk, and the market's equity volatility overshoots it.}{E3, M4}
\crow{C3}{Guarding the floor that regulation makes binding, with the movement of book capital as the risk, beats the standard model on banks and on a year it had not seen, without placing the level of buffers on that year.}{E4, E5, E6, E7, M3, L7, L8}
\crow{C4}{The floor the enhanced model guards rests on sources rather than on fit, the well-capitalised line before 2010 and the minimum plus the conservation buffer after.}{M5, L7, L8, L9, L14}
\crow{C6}{Observed buffers imply a risk about three times earnings volatility on 2005 to 2014, against Peura and Keppo's 1.1 to 1.3 times on their 1983 to 2002 subsample (L15), with loss-based measures at about one and the market at about eight. Juselius, Ponte Marques and Tarashev's buffer of about four standard deviations of buffer shocks has a different denominator (L19), and Peura and Jokivuolle invert buffers for a confidence level rather than a risk scale (L23). No paper read measures this risk scale against loss-based and market measures.}{E2, E3, L13, L15, L18, L19, L23}
\bottomrule
\end{longtable}
\end{center}

\begin{center}\footnotesize
\begin{longtable}{P{0.5cm}P{13.0cm}P{2.6cm}}
\caption{Pending.}\label{tab:pending}\\
\toprule ID & Statement & Rows \\ \midrule \endfirsthead
\toprule ID & Statement & Rows \\ \midrule \endhead
\crow{C5}{Earlier work fits a single model to the data it then explains, Peura and Keppo by calibration and Corbae and D'Erasmo by estimation to buffer levels in sample (L1, L20), and Dai, Huang and Keppo compare two versions of one model out of sample on other targets (L17). No paper read scores the published dynamic models of different authors against each other on pre-registered buffer facts out of sample, with no parameter fitted to those facts. \pend{Migueis (2017), a recalibration of Peura and Keppo, seen in abstract only.}}{L1, L17, L20}
\crow{C7}{The size gradient of buffers survives controls for access to capital by affiliation and listing, the explanation Peura and Keppo give for it (L16). Sengupta and Hogue find banks without a holding company hold more within banks on another definition and period (L22), and Corbae and D'Erasmo assign the gradient to funding volatility and issuance costs without testing either (L21). \pend{a test of funding volatility, and a split of stand-alone banks into those with no holding company and those of one-bank holding companies, each under its own pre-registration.}}{E9, L6, L16, L21, L22}
\bottomrule
\end{longtable}
\end{center}

\bibliographystyle{plainnat}
\bibliography{../refs}

\appendix

\section{Proofs}\label{app:proofs}

\paragraph{Notation.} $x$ is the buffer over the floor, $\fl$ the floor,
$R=0.105$ the reference, $\mu_i$ and $\sigma_i$ the mean and the
volatility of bank $i$'s yearly earnings, $s_i$ the judged risk, $\rho$ the
discount rate, $p=d_+(\mu,s,\rho)$ and $m=d_-(\mu,s,\rho)$ the roots,
$a_1=m/(pm-p^2)$, $a_2=p/(pm-m^2)$, $u_0$ the dividend barrier, $D$ a
decline in the capital ratio and $z$ a normal quantile.

\begin{mproof}{M1}{PeuraKeppo.char_eq, PeuraKeppo.dplus_pos, PeuraKeppo.dminus_neg}
\[ d_\pm=\frac{-\mu\pm\sqrt{\mu^2+2\rho s^2}}{s^2},\qquad \tfrac12 s^2d_\pm^2+\mu d_\pm-\rho=0. \]
With $\rho>0$, $\sqrt{\mu^2+2\rho s^2}>|\mu|$, so $d_+>0$ and $d_-<0$.
\end{mproof}

\begin{mproof}{M2}{PeuraKeppo.no_issue_barrier, PeuraKeppo.coeff_signs, MilneWhalley.buffer_eq_u0}
\[ e^{(p-m)u_0}=\Big(\frac{-m}{p}\Big)^{2},\quad a_1e^{-pu_0}+a_2e^{-mu_0}=e^{-pu_0}\Big(a_1+a_2\frac{m^2}{p^2}\Big)=0, \]
\[ a_1=\frac{m}{p(m-p)}>0>a_2=\frac{p}{m(p-m)},\qquad \frac{1}{p-m}\log\frac{m^2}{p^2}=\frac{2}{p-m}\log\frac{-m}{p}. \]
\end{mproof}

\begin{mproof}{M3}{MilneWhalley.buffer_shift_invariant}
\[ (C_t+d)-(C_0+d)=C_t-C_0, \]
so with $k_t=\fl+x_t$ a change $\fl\mapsto\fl+d$ moves $k_t$ by $d$ and leaves $x_t-x_0$, $u_0$ and the closure time $\inf\{t:x_t=0\}$ unchanged.
\end{mproof}

\begin{mproof}{M4}{StressTests.judged_sigma, StressTests.sqrt_horizon, StressTests.decline_of_rows}
\[ \sqrt{9/4}=\tfrac32,\qquad \frac{D}{z\sqrt{9/4}}=\frac{D}{\tfrac32 z},\qquad m\le a\ \Rightarrow\ D=a-m\ge0. \]
\end{mproof}

\begin{mproof}{M5}{OrozcoRubio.basel3_floor_eq_R, OrozcoRubio.basel3_floor_above_pca}
\[ 0.08+0.025=0.105=R,\qquad 0.10<0.08+0.025. \]
\end{mproof}

\section{Evidence and verification}\label{app:evidence}

The Lean theorems are those of \path{lean/mathlib/}, built and audited by the
parent bundle (\path{verify.sh}). The checker
\path{enhancement/checks/verify_enhancement.py} applies MS-1 to MS-9 and
re-runs every results file an empirical row cites.

\paragraph{What Lean does not carry.} The Brownian motion of A1, the
optimality of paying at $u_0$ (A2), and the simulation of the facts. The
facts are computed on pinned data and reproduced exactly by the checker,
which is the third class of evidence of this bundle.

\paragraph{Illustrations.} No row rests on SymPy or sampling outside a
pinned, reproduced results file.

\section{Terminology and measurement map}\label{app:terms}

The map names every object of the model once, and the manuscript uses these
names and no others. Inputs are what the model takes as given, internal
objects are built from the inputs inside the model, and outputs are what the
rows predict. Directly observed is Yes when a record of the world gives the
object without computation, which is when the last column begins with
Recorded. Every computed object names its recovery formula and the Lean
theorem that carries it. The same tables make the standalone
\path{enhancement/MEASUREMENT_MAP.tex}.

\subsection{Inputs}

\footnotesize
\begin{longtable}{@{}P{1.4cm}P{2.3cm}P{2.2cm}P{3.6cm}P{1.2cm}P{4.5cm}@{}}
\caption{Inputs of the model.}\label{tab:inputs}\\
\toprule
\textbf{Symbol} & \textbf{Name} & \textbf{Type} & \textbf{Represents} & \textbf{Directly observed} & \textbf{Obtained how} \\
\midrule
\endfirsthead
\toprule
\textbf{Symbol} & \textbf{Name} & \textbf{Type} & \textbf{Represents} & \textbf{Directly observed} & \textbf{Obtained how} \\
\midrule
\endhead
\bottomrule
\endfoot
\mm{$k$}{capital ratio}{A real number above zero.}{Total capital, tier 1 plus tier 2, over risk-weighted assets, for each bank and quarter.}{No}{Computed as \path{(RBCT1J+RBCT2)/RWAJT} from the Call Report, fields defined in the record lit/fdic\_2026. The ratio is free of the unit of account (\lean{EnhancementMap.capital_ratio_units}).}
\mm{$\fl$}{regulatory floor}{A real number between zero and one.}{The capital ratio at which the bank is closed in the model (A3, A5).}{No}{Assumed by A5 at the lines the sources find binding, 8\% for the baseline, and 10\% to 2009Q4 and 10.5\% from 2010Q1 for the enhanced model (M5).}
\mm{$R$}{reference}{A real number, 0.105.}{The ratio the facts are measured from.}{No}{Assumed, fixed in comparison/PREREG\_comparison.md before any fact was computed.}
\mm{$\mu_i$}{mean earnings}{A real number.}{The mean of the bank's yearly net income over the previous year's risk-weighted assets, less the riskless return on its capital ratio (A1).}{No}{Computed over ten years of December reports by the code of replication/peura\_keppo\_2006, each year as net income over risk-weighted assets less the riskless rate times the capital ratio (\lean{EnhancementMap.net_earnings}).}
\mm{$\sigma_i$}{earnings volatility}{A positive real number.}{The standard deviation of the same yearly earnings (A6).}{No}{Computed over the same ten years from the same yearly terms (\lean{EnhancementMap.net_earnings}).}
\mm{$\rho$}{discount rate}{A positive real number, 0.04.}{The rate at which the bank discounts dividends, its cost of equity (A4).}{No}{Assumed at Peura and Keppo's value.}
\mm{$s_i$}{judged risk}{A positive real number.}{The volatility the buffer over the floor moves with in the model (A1, A6).}{No}{Assumed by A6 to be the earnings volatility or the book-capital volatility, or one of the measures of E3.}
\mm{$\sigma_{D,i}$}{book-capital volatility}{A positive real number.}{The standard deviation of the yearly change in the capital ratio over ten years (A6).}{No}{Computed by comparison/book\_vol.py. The change is the same for the capital ratio and for the buffer over the floor (\lean{EnhancementMap.change_free_of_floor}).}
\mm{$\sigma_{E,i}$}{equity volatility}{A positive real number.}{The annualised standard deviation of daily log changes of the parent's share price.}{No}{Computed from pinned daily prices by comparison/judged\_risk.py as $\sqrt{252}$ times the daily standard deviation (\lean{EnhancementMap.annualise}).}
\mm{$X_i$}{mean capital ratio}{A real number above zero.}{The mean of the bank's capital ratio over ten December reports.}{No}{Computed by comparison/judged\_risk.py; a mean lies between the least and the largest ratio (\lean{EnhancementMap.mean_between}).}
\mm{$\sigma_{B,i}$}{market judged risk}{A positive real number.}{The volatility of the capital ratio if book capital moved with the market value of equity (E3).}{No}{Computed as $\sigma_{E,i}X_i$, the standard deviation of $X_i$ times a return (\lean{EnhancementMap.market_scale}).}
\mm{$D$}{stress-test decline}{A real number at least zero.}{The fall in a holding company's tier 1 common ratio from its actual value to its projected minimum under the severely adverse scenario.}{No}{Computed from the published rows of lit/dfast\_2013 and lit/dfast\_2014 as the actual less the minimum (\lean{StressTests.decline_of_rows}).}
\mm{$z$}{normal quantile}{A positive real number.}{The quantile of a normal loss that a decline is read as (M4).}{No}{Assumed on the grid 1.645, 2.326, 3.090, with 2.326 central.}
\mm{$\sigma_{A,i}$}{stress-test judged risk}{A positive real number.}{The volatility under which a nine-quarter decline $D$ is the $z$-quantile (E3).}{No}{Computed as $D/(\tfrac32 z)$ (\lean{StressTests.judged_sigma}).}
\mm{$\sigma_{C,i}$}{own-loss judged risk}{A positive real number.}{The same conversion applied to the bank's own worst nine-quarter decline (E3).}{No}{Computed by comparison/own\_loss.py as $D/(\tfrac32 z)$ with the bank's decline (\lean{StressTests.judged_sigma}).}
\mm{$\sigma_{N,i}$}{do-nothing volatility}{A positive real number.}{The standard deviation of the yearly change in Berger et al.'s do-nothing ratio (E8).}{No}{Computed by comparison/do\_nothing.py; each change is an earnings part plus an asset part (\lean{EnhancementMap.do_nothing_split}).}
\mm{$S_i$}{bank size}{A real number.}{The log of the bank's mean total assets (E9).}{No}{Computed from the Call Report field \path{ASSET}; the log preserves the order of sizes (\lean{EnhancementMap.size_monotone}).}
\end{longtable}

\subsection{Internal objects}

\footnotesize
\begin{longtable}{@{}P{1.4cm}P{2.3cm}P{2.2cm}P{3.6cm}P{1.2cm}P{4.5cm}@{}}
\caption{Internal objects of the model.}\label{tab:internal}\\
\toprule
\textbf{Symbol} & \textbf{Name} & \textbf{Type} & \textbf{Represents} & \textbf{Directly observed} & \textbf{Obtained how} \\
\midrule
\endfirsthead
\toprule
\textbf{Symbol} & \textbf{Name} & \textbf{Type} & \textbf{Represents} & \textbf{Directly observed} & \textbf{Obtained how} \\
\midrule
\endhead
\bottomrule
\endfoot
\mm{$x$}{buffer over the floor}{A real number, the state.}{The capital ratio less the regulatory floor (A1, A3).}{No}{Computed as $k-\fl$ (\lean{EnhancementMap.buffer_of_ratio}).}
\mm{$W_t$}{Brownian motion}{A standard Brownian motion.}{The source of the buffer's random movement (A1).}{No}{Never observed.}
\mm{$p$}{positive root}{A positive real number.}{The positive root of the value equation (M1).}{No}{Computed from $\mu_i$, $s_i$ and $\rho$ (\lean{PeuraKeppo.dplus_pos}).}
\mm{$m$}{negative root}{A negative real number.}{The negative root of the value equation (M1).}{No}{Computed from $\mu_i$, $s_i$ and $\rho$ (\lean{PeuraKeppo.dminus_neg}).}
\mm{$a_1$, $a_2$}{barrier coefficients}{Two real numbers of opposite sign.}{The weights of the two exponentials in the smooth-fit condition (M2).}{No}{A device of the proof of M2 (\lean{PeuraKeppo.coeff_signs}).}
\mm{$u_0$}{dividend barrier}{A positive real number.}{The buffer over the floor at which dividends are paid (A2, M2).}{No}{Computed as $\frac{2}{p-m}\log\frac{-m}{p}$ (\lean{PeuraKeppo.no_issue_barrier}).}
\end{longtable}

\subsection{Outputs}

\footnotesize
\begin{longtable}{@{}P{1.4cm}P{2.3cm}P{2.2cm}P{3.6cm}P{1.2cm}P{4.5cm}@{}}
\caption{Outputs of the model.}\label{tab:outputs}\\
\toprule
\textbf{Symbol} & \textbf{Name} & \textbf{Type} & \textbf{Represents} & \textbf{Directly observed} & \textbf{Obtained how} \\
\midrule
\endfirsthead
\toprule
\textbf{Symbol} & \textbf{Name} & \textbf{Type} & \textbf{Represents} & \textbf{Directly observed} & \textbf{Obtained how} \\
\midrule
\endhead
\bottomrule
\endfoot
\mm{$F_1$}{thin-buffer dividend share}{A share between zero and one.}{The share of quarters with dividends paid from a buffer over the reference below $\sigma_i$, observed or predicted.}{No}{Computed by comparison/facts.py and window.py on the data, and by the simulations on the model (\lean{EnhancementMap.share_unit}).}
\mm{$F_2$}{far-above share}{A share between zero and one.}{The share of bank-quarters more than $4\sigma_i$ above the reference, observed or predicted.}{No}{Computed by the same scripts (\lean{EnhancementMap.share_unit}).}
\end{longtable}
\normalsize


\section{Primitives and scope}\label{app:scope}

\paragraph{Process.} A1 is Peura and Keppo's process for capital over the
minimum, with the bank's own mean and volatility. M1 and M2 depend on it.

\paragraph{Issuance.} A2 drops issuance. Milne and Whalley show the barrier
without issuance bounds the buffer with costly issuance from above (L3), so
the enhanced model's buffer is the largest the family gives. M2 depends on it.

\paragraph{Closure.} A3 closes the bank at the floor. The facts are measured
from $R$, not from the floor, so the floor moves the barrier in the ratio
(M3) without changing the facts' yardstick.

\paragraph{Discounting.} A4 takes Peura and Keppo's cost of equity of four
percent. The training-window scale check (E2) reports lower values.

\paragraph{Floor.} A5 is the author's choice of 25 October 2026, taken from
Orozco and Rubio (L7, L8, L9) and Lang and Menno (L14). It was made after a
floor near $R$ had been seen to work on 2015 to 2024, which is why the 2025
hold-out (E6, E7) decides the frozen model.

\paragraph{Judged risk.} A6 replaces the earnings volatility by the
book-capital volatility, which includes the bank's own capital actions (E8).
The alternatives scored and rejected, the stress-test decline, the own worst
decline and the equity volatility, are E3. Converting a decline to a volatility
assumes a normal quantile $z$, and a stress scenario has no probability of
its own (L24). Credit losses are long-tailed, so a normal quantile if anything
overstates the volatility a decline implies, and the stress-test judged risk
of E3 is an upper value.

\end{document}
